\providecommand{\ggd}{\mathfrak{g}_{\Delta_5}}
\providecommand{\Lambdaa}{\Lambda^{3,2}}
\providecommand{\Lambdab}{\Lambda^{2,1}}
\providecommand{\Lambdac}{\Lambda^{2,1}_{II}}
\providecommand{\WII}{W^{(2)}(\Lambdac)}
\providecommand{\PII}{\mathcal{P}_{II}}
\providecommand{\PIIprim}{\mathcal{P}_{II,\mathrm{prim}}}
\providecommand{\HIV}{\mathcal{H}^{IV}_{+}}
\providecommand{\Cone}{\mathcal{C}}

\section*{Root Datum, Cartan Matrix, and Denominator of $\ggd$}
\label{sec:delta5-root-cartan}

Lorgat~2020, following Borcherds and Gritsenko--Nikulin, constructs the
generalised Borcherds--Kac--Moody superalgebra $\ggd$ from the Borcherds
lift of the weight-$0$ index-$1$ weak Jacobi form $\phi_{0,1}$ to the
Igusa cusp form $\Delta_5$.  The construction uses the signature-$(3,2)$
lattice $\Lambdaa$, the type-II hyperbolic sublattice $\Lambdac$, the
reflection group $\WII$, the Weyl vector $\rho$, and the three primitive
real simple roots
\[
\PIIprim=\{\delta_1,\delta_2,\delta_3\}.
\]
The Macdonald--Weyl--Kac--Borcherds denominator identity of $\ggd$ is
the product expansion of $64^{-1}\Delta_5(2Z)$ at the zero-dimensional
cusp.

\subsection*{The Pfaffian Lattice $\Lambdaa$}

Let $\Lambda_4=\mathbb{Z}e_1\oplus\mathbb{Z}e_2\oplus
\mathbb{Z}e_3\oplus\mathbb{Z}e_4$, and define the Pfaffian pairing on
$\wedge^2\Lambda_4$ by
\[
u\wedge v=(u,v)\,e_1\wedge e_2\wedge e_3\wedge e_4 .
\]
For
\[
q=e_1\wedge e_3+e_2\wedge e_4,
\qquad
\Lambdaa=q^\perp\subset\wedge^2\Lambda_4,
\]
the basis
\begin{equation}
f_1=e_1\wedge e_2,\quad
f_2=e_2\wedge e_3,\quad
f_3=e_1\wedge e_3-e_2\wedge e_4,\quad
f_{-2}=e_4\wedge e_1,\quad
f_{-1}=e_4\wedge e_3
\label{eq:delta5-pfaffian-basis}
\end{equation}
has Gram matrix
\begin{equation}
G(\Lambdaa)=
\begin{pmatrix}
0&0&0&0&-1\\
0&0&0&-1&0\\
0&0&2&0&0\\
0&-1&0&0&0\\
-1&0&0&0&0
\end{pmatrix}.
\label{eq:delta5-gram-lambda32}
\end{equation}
The two hyperbolic pairs are $(f_1,f_{-1})$ and $(f_2,f_{-2})$; the
orthogonal vector $f_3$ has square $2$.  Thus
\[
\Lambdaa\simeq U\oplus U\oplus\langle 2\rangle,
\qquad
\operatorname{sign}(\Lambdaa)=(3,2),
\qquad
\det G(\Lambdaa)=2.
\]
The discriminant group is $\Lambdaa^\vee/\Lambdaa\simeq\mathbb{Z}/2\mathbb{Z}$,
generated by the $\langle2\rangle$ summand, with discriminant value
$1/2$ modulo $2\mathbb{Z}$.

The Pfaffian form is $\mathrm{SL}_4(\mathbb{Z})$-invariant.  The stabiliser
of $q$ is $\mathrm{Sp}_4(\mathbb{Z})$, and the induced exterior-square
representation gives
\[
\wedge^2:\mathrm{Sp}_4(\mathbb{Z})/\{\pm I_4\}
\xrightarrow{\sim}\mathrm{SO}^{+}(\Lambdaa).
\]
The type-IV domain $\HIV$ attached to $\Lambdaa$ is thereby identified
with the Siegel upper half-space $\mathbb{H}_2$.

\subsection*{The Type-II Hyperbolic Sublattice}

The primitive hyperbolic sublattice used for the denominator chamber is
\[
\Lambdab=\mathbb{Z}f_2\oplus\mathbb{Z}f_3\oplus\mathbb{Z}f_{-2}
\subset\Lambdaa,
\]
with Gram matrix
\[
\begin{pmatrix}
0&0&-1\\
0&2&0\\
-1&0&0
\end{pmatrix}
\]
in the ordered basis $(f_2,f_3,f_{-2})$.  Its type-II reflection sublattice is
\[
\Lambdac=\{m f_2+\ell f_3+n f_{-2}\in\Lambdab\mid m\equiv n\equiv0
\pmod 2\}.
\]
The three primitive type-II real simple roots are
\begin{equation}
\delta_1=2f_2-f_3,\qquad
\delta_2=2f_{-2}-f_3,\qquad
\delta_3=f_3 .
\label{eq:delta5-real-simple-roots}
\end{equation}
They form a basis of $\Lambdac$.  Their Gram matrix is
\begin{equation}
G_{\mathrm{BKM}}=((\delta_i,\delta_j))_{i,j=1}^{3}=
\begin{pmatrix}
2&-2&-2\\
-2&2&-2\\
-2&-2&2
\end{pmatrix}.
\label{eq:delta5-cartan}
\end{equation}
Since each real simple root has square $2$, the symmetric Gram matrix
\eqref{eq:delta5-cartan} is also the generalised Cartan matrix:
\[
A_{ij}=\frac{2(\delta_i,\delta_j)}{(\delta_j,\delta_j)}
=(\delta_i,\delta_j).
\]
Its determinant is $-32$ and its eigenvalues are $4,4,-2$, so the Cartan
space has signature $(2,1)$.  The Cartan rank of $\ggd$ is therefore $3$.

For $i\neq j$,
\[
\frac{A_{ij}A_{ji}}{A_{ii}A_{jj}}=1,
\]
so the Coxeter exponent is $m_{ij}=\infty$.  The Coxeter chamber
$\PII\subset\Cone(\Lambdab)_+/\mathbb{R}_{>0}$ is an ideal hyperbolic
triangle.  Its three ideal vertices are represented by the primitive
null vectors
\[
2f_2,\qquad 2f_{-2},\qquad 2f_2-2f_3+2f_{-2}.
\]
The reflection group $\WII$ is generated by the reflections in
$\delta_1,\delta_2,\delta_3$, and $\mathrm{Aut}(\PII)\simeq S_3$ permutes
the three ideal vertices.

The real simple roots of $\ggd$ are exactly the three roots in
\eqref{eq:delta5-real-simple-roots}.  The full positive real-root set is
\begin{equation}
\Delta^{\mathrm{re}}_+(\ggd)=
\WII\cdot\{\delta_1,\delta_2,\delta_3\}.
\label{eq:delta5-real-root-orbit}
\end{equation}
The integer $10$ in the product formula
$\Delta_5=\prod\nu_{a,b}$ counts the ten even genus-$2$ theta constants;
it is unrelated to the number of simple real roots.

\subsection*{The Weyl Vector}

A lattice Weyl vector for the real chamber satisfies
\[
(\rho,\delta_i)=-\frac{(\delta_i,\delta_i)}2=-1,\qquad i=1,2,3.
\]
Writing $\rho=c_1\delta_1+c_2\delta_2+c_3\delta_3$, this is the linear
system
\[
G_{\mathrm{BKM}}
\begin{pmatrix}c_1\\c_2\\c_3\end{pmatrix}
=
\begin{pmatrix}-1\\-1\\-1\end{pmatrix}.
\]
Since $\det G_{\mathrm{BKM}}=-32$, the solution is unique.  Subtracting
any two equations gives $c_1=c_2=c_3$, and then
$2c-4c=-1$, hence
\begin{equation}
\rho=\frac12(\delta_1+\delta_2+\delta_3)
=f_2-\frac12 f_3+f_{-2}.
\label{eq:delta5-weyl-vector}
\end{equation}
With the Gram matrix \eqref{eq:delta5-gram-lambda32},
\[
(\rho,\rho)=2(f_2,f_{-2})+\frac14(f_3,f_3)
=-2+\frac12=-\frac32.
\]
The positive cone is taken with the convention
\[
\Cone(\Lambdab)_+=\{x\in\Lambdab\otimes\mathbb{R}\mid (x,x)<0\}.
\]
Thus $\rho$ lies in the interior of the positive cone and in
$\Lambdac^\vee\cap\tfrac12\Lambdac$, with
$(\rho,\delta_i)=-1$ for all real walls.

\subsection*{Jacobi Coefficients and Imaginary Simple Roots}

Let
\[
\phi_{0,1}(\tau,z)=\sum_{n\geq0,\ell\in\mathbb{Z}}
f(n,\ell)\,q^n r^\ell
\]
be the Eichler--Zagier weight-$0$ index-$1$ weak Jacobi form, normalised by
\[
f(0,0)=10,\qquad f(0,-1)=f(0,1)=1,
\]
and
\[
\phi_{0,1}(\tau,z)=
(r^{-1}+10+r)
q(10r^{-2}-64r^{-1}+108-64r+10r^2)+\cdots .
\]
The coefficient $f(n,\ell)$ depends only on the discriminant
$4n-\ell^2$.

The Borcherds product is
\begin{equation}
\frac1{64}\Delta_5(Z)
=
e^{\pi i(z_1+z_2+z_3)}
\prod_{\substack{(n,\ell,m)>0}}
\left(1-e^{2\pi i(nz_1+\ell z_2+mz_3)}\right)^{f(nm,\ell)},
\label{eq:delta5-borcherds-product}
\end{equation}
where $(n,\ell,m)>0$ means $n,m\geq0$, $\ell\in\mathbb{Z}$ if
$n>0$ or $m>0$, and $\ell<0$ if $n=m=0$.  The denominator function of
$\ggd$ is
\begin{equation}
\Phi(z)=
e^{-2\pi i(\rho,z)}
\prod_{\alpha\in\Delta_+}
\left(1-e^{-2\pi i(\alpha,z)}\right)^{\mathrm{mult}(\alpha)}
=\frac1{64}\Delta_5(2Z).
\label{eq:delta5-denominator-product}
\end{equation}
The exponent $f(nm,\ell)$ in \eqref{eq:delta5-borcherds-product} is the
superdimension
\[
\mathrm{mult}(\alpha)=
\dim\ggd_{\alpha,\bar{0}}-\dim\ggd_{\alpha,\bar{1}}
\]
of the corresponding positive root sector in \eqref{eq:delta5-denominator-product}.
The three real simple roots have multiplicity $1$, read from
$f(0,-1)=1$ in the $2Z$ normalisation.

The imaginary simple roots are indexed by
$a\in\Lambdac\cap\mathbb{R}_{>0}\PII$ with $(a,a)\leq0$.  The even
lightlike roots are
\begin{equation}
\Delta^{\mathrm{im}}_{\bar0}
=\{\tau(a)a\mid (a,a)=0,\ \tau(a)>0\},
\qquad
\tau(a)=9.
\label{eq:delta5-even-imaginary-roots}
\end{equation}
For a primitive null vertex $a_0$, the equality
\[
1+\frac1{64}\sum_{t\geq1} f(1+2t,1,1)q^t
=\prod_{k\geq1}(1-q^k)^9
\]
gives $\tau(ta_0)=9$ for all $t\geq1$, and the $S_3$-symmetry of
$\PII$ transfers this value to the other two null vertices.  The odd
imaginary simple roots are governed by the same Fourier-transform
coefficients $m(a)$:
\begin{equation}
\Delta^{\mathrm{im}}_{\bar1}
=\{m(a)a\mid (a,a)<0,\ m(a)<0\}.
\label{eq:delta5-odd-imaginary-roots}
\end{equation}
Negative Jacobi coefficients give odd root spaces; positive coefficients
give even root-space superdimension.  For example,
\[
f(1,0)=108,\qquad f(1,\pm1)=-64,\qquad f(1,\pm2)=10.
\]
The real-root Serre relations alone therefore define only the rank-$3$
Kac--Moody core.  The Borcherds superalgebra $\ggd$ also includes the
imaginary simple roots \eqref{eq:delta5-even-imaginary-roots} and
\eqref{eq:delta5-odd-imaginary-roots}, with multiplicities fixed by
$\phi_{0,1}$.

\subsection*{The Denominator Identity}

The Weyl--Kac--Borcherds character formula applied to the trivial
one-dimensional representation gives
\begin{equation}
\Phi(z)=
\sum_{w\in\WII}\det(w)
\left(
e^{-2\pi i(w\rho,z)}
-
\sum_{a\in\Lambdac\cap\mathbb{R}_{>0}\PII}
m(a)e^{-2\pi i(w(\rho+a),z)}
\right),
\label{eq:delta5-weyl-sum}
\end{equation}
and the product expression is \eqref{eq:delta5-denominator-product}.
Equations \eqref{eq:delta5-weyl-sum} and
\eqref{eq:delta5-denominator-product} are the Macdonald identity for
$\ggd$:
\[
\sum_{w\in\WII}\det(w)\,
w\left(
e^\rho
\prod_{\alpha\in\Delta^{\mathrm{im}}_+}
(1-e^{-\alpha})^{\mathrm{mult}(\alpha)}
\right)
=
e^\rho
\prod_{\alpha\in\Delta_+}(1-e^{-\alpha})^{\mathrm{mult}(\alpha)}.
\]
The equality with $64^{-1}\Delta_5(2Z)$ is the
Borcherds--Gritsenko--Nikulin denominator theorem in the normalisation
used above.

\subsection*{K3$\times E$ Scope and the Four Invariants}

The BKM denominator branch over $K3\times E$ is the Borcherds product
branch governed by $\ggd$ and by the compact Hall object under the compact
Hall hypotheses:
\[
D^\wedge_\hbar\!\left(Y^+_{\mathrm{Hall}}(K3\times E)\right).
\]
This branch is separate from the historical K3 self-mirror Yangian branch,
which quantises the rank-$24$ Mukai--Heisenberg data of K3.  The common
geometric background leaves the two quantum groups distinct.

The four invariants attached to the $K3\times E$ example live on distinct
constructions:
\[
\kappa_{\mathrm{cat}}(K3\times E)=0,\qquad
\kappa_{\mathrm{ch}}^{\mathrm{Heis}}(K3\times E)=3,\qquad
\kappa_{\mathrm{BKM}}(\Delta_5)=5,\qquad
\kappa_{\mathrm{fiber}}(K3\times E)=24.
\]
Here $\kappa_{\mathrm{cat}}$ is $\chi(\mathcal{O}_{K3})\chi(\mathcal{O}_E)=2\cdot0$,
$\kappa_{\mathrm{BKM}}(\Delta_5)$ is the Borcherds weight, and
$\kappa_{\mathrm{fiber}}$ is the Mukai-lattice rank
\[
\operatorname{rk}\widetilde H(K3,\mathbb{Z})=1+22+1=24.
\]
The Cartan rank of $\ggd$ is $3$, the set of Borcherds imaginary simple
roots is infinite, and the Mukai lattice has rank $24$.  These are three
different counts.

For the primitive CHL denominator ladder
$N\in\{1,2,3,4,6\}$, the Borcherds weight formula is
\[
\kappa_{\mathrm{BKM}}(\Phi_N)=\frac{c_N(0)}2,
\]
with constants
\[
\begin{array}{c|ccccc}
N&1&2&3&4&6\\
\hline
c_N(0)&10&8&6&4&2\\
\kappa_{\mathrm{BKM}}(\Phi_N)&5&4&3&2&1
\end{array}
\]
The $N=1$ entry is $\Delta_5$: $c_1(0)=10$ and
$\kappa_{\mathrm{BKM}}(\Delta_5)=5$.

The toric comparison uses the positive half
\[
\mathrm{CoHA}(\mathbb{C}^3)=Y^+(\widehat{\mathfrak{gl}}_1).
\]
The full $\mathcal{W}_{1+\infty}$ object appears after the appropriate
Drinfeld double, centre, completion, or evaluation passage.  This is the
same positive-half versus doubled-object distinction required in the
$K3\times E$ Hall--Drinfeld comparison.
